%%
%%                  TEMPLATE for Math 402/436 project report
%%
\documentclass[11pt]{amsart}
%%% WARNING: Do NOT change the page size, fonts, or margins!  Penalties will apply.


\usepackage{graphicx}
\usepackage{amssymb,amsmath,amsthm}
\usepackage{placeins} %enables \FloatBarrier, which prevents figures and tables from going below it.
\usepackage{hyperref} %makes cross references into hyperlinks. 
\usepackage{caption}
\usepackage{subcaption} %Allows subfigures
 

%%% WARNING: Do NOT change the page size, fonts, or margins!  Penalties will apply.
%%% WARNING: Do NOT change the page size, fonts, or margins!  Penalties will apply.

\begin{document}

\title{Modeling the Dynamic Conditional Information Coefficient}
\author{Conner B. Draper, Seokkwon Yoon}

% Change the date to match the date you actually wrote this paper
\date{\today} 

\maketitle % this actually makes the title

\begin{abstract}
The Information Coefficient (IC) is a widely used measure of predictive signal quality in quantitative finance, typically treated as a constant. However, predictive power may depend on signal strength and vary over time. In this paper, we study a dynamic conditional Information Coefficient, denoted $IC_t(z)$, where $z$ represents standardized signal strength. We compare five approaches: a constant IC baseline, independent binned Kalman filters, a Kalman filter polynomial state-space model, Recursive Least Squares with RBFs, and Nadaraya-Watson kernel regression. We evaluate their effectiveness across eight cross-sectional signals through walk-forward mean-variance optimization, determining whether modeling conditional predictive strength leads to improved risk-adjusted performance.
\end{abstract}

%% First Section
\section{Research question and overview of the data }

In traditional quantitative finance, the Information Coefficient (IC) is often assumed to be constant (e.g., $IC \approx 0.05$). This simplification ignores the possibility that predictive power depends on signal strength and market regimes. Our central research question is: \textit{Does predictive performance vary as a function of standardized signal values $z$, and can online learning models dynamically estimate this relationship to improve portfolio allocation?}

The central object of interest is $IC_t(z)$, representing a conditional and time-varying Information Coefficient. If signals with larger magnitudes (high $|z|$) exhibit stronger predictive power than those near zero, portfolio construction methods that exploit this structure may significantly outperform static approaches.

We consider a panel dataset consisting of standardized signals $z_{i,t}$ and subsequent realized returns $r_{i,t+1}$ for a cross-section of assets over time. We evaluate eight signals: style, industry, and idiosyncratic momentum; style, industry, and idiosyncratic reversal; idiosyncratic volatility; and betting against beta (BAB). 

\begin{itemize}
    \item Data Weaknesses: Financial return data is notoriously noisy, and the assumption of a constant variance structure is often violated. Furthermore, micro-cap stocks can introduce extreme outliers.
    \item Suitability: This dataset is highly suitable because it spans diverse risk factors (Barra components and idiosyncratic metrics), providing the breadth required by the Fundamental Law of Active Management \cite{Vandermeersch} to validate statistical learning models.
    \item Expected Analysis: We expect that models capturing time-varying, state-dependent predictive power will yield superior expected returns ($\mu_t$) and smoothly down-weight decaying signals compared to the constant baseline.
\end{itemize}

\begin{figure}[tph] %Put your images in a figure like this
     \centering
      \begin{subfigure}[b]{0.45\textwidth}
         \centering
	   \includegraphics[width=.95\linewidth]{crossover1.png}
     \end{subfigure}
     \begin{subfigure}[b]{0.45\textwidth}
         \centering
	   \includegraphics[width=.95\linewidth]{crossover2.png}
     \end{subfigure}
     \caption[Behavior of Conditional Returns]{In these plots we consider the behavior of the estimated conditional return curves. The left plot illustrates the learned piece-wise relationship using the binning method across the z-score distribution. The right plot shows the time-varying conditional return curve estimated via Kalman filtering. The dynamic models successfully adapt to non-linearities at the tails of the distribution, which a static baseline ignores.}
\label{fig:MeanSquaredError} 
\end{figure}

%% Second Section 
\section{Data Cleaning / Feature Engineering }

To avoid look-ahead bias, we ensure that all signals at time $t$ are used only to predict returns at time $t+1$. The dataset is split into chronological training and testing periods. 

The raw signals are engineered by standardizing them cross-sectionally on each date, yielding $z_{i,t} \in [-3, 3]$. Furthermore, to properly model the IC using the Grinold framework, we risk-normalize the target variable. Instead of raw returns, the target is $y_{i,t+1} = r_{i,t+1} / \sigma_{i,t}$, where $\sigma_{i,t}$ is the specific risk. We apply a minimum price filter ($>\$5.00$) to filter out illiquid micro-cap noise. This transforms the modeling objective into predicting risk-adjusted returns where the scaling factor directly represents the IC.

%% Third Section
\section{Data Visualization and Basic Analysis}

Summary statistics of the rolling signals indicate that predictive power fluctuates widely. Scatter plots of $y_{i,t+1}$ versus $z_{i,t}$ reveal a high noise-to-signal ratio, but binning the cross-sectional data by $z$-score deciles shows observable, often non-linear trends at the extremes. 

By analyzing the signal-return relationship over a rolling 120-day window, we observe that the empirical IC frequently deviates from 0.05, sometimes flipping signs entirely during momentum crashes. Using visualizations similar to Figure~\ref{fig:MeanSquaredError}, we identified that the signal strength is asymmetric; for example, high positive values of idiosyncratic momentum often show a different predictive slope than deeply negative values.

%%Fourth Section
\section{Learning Algorithms and In-depth Analysis }

We model predictive strength utilizing five approaches. Under the standard Grinold formula, the expected return is $\mu_{i,t} = \sigma_{i,t} \cdot IC_t(z_{i,t}) \cdot z_{i,t}$.

\subsection{Static Baseline}
The traditional approach assumes a constant Information Coefficient: $IC = 0.05$.

\subsection{Binned Kalman Method}
We partition the signal space into bins $\mathcal{B}_1, \dots, \mathcal{B}_{15}$ and utilize 15 independent 1-D Kalman filters. This produces a piecewise estimate of the conditional expected risk-normalized return.

\subsection{Kalman Filter Polynomial Model}
To capture continuous time-varying relationships, we model the IC as a polynomial function of $z$:
\[
IC_t(z) = a_t + b_t z + c_t z^2.
\]
We define the state vector $\theta_t = \begin{bmatrix} a_t & b_t & c_t \end{bmatrix}^T$. The state evolution follows a random walk:
\[
\theta_t = \theta_{t-1} + w_t, \quad w_t \sim \mathcal{N}(0,Q).
\]
Since our target is $y_{i,t+1} = IC_t(z_{i,t}) \cdot z_{i,t}$, the observation equation becomes:
\[
y_{i,t+1} = \begin{bmatrix} z_{i,t} & z_{i,t}^2 & z_{i,t}^3 \end{bmatrix} \theta_t + \epsilon_{i,t}.
\]
The Kalman filter recursively updates $\theta_t$ daily, allowing the signal-return relationship to evolve over time.

\subsection{Non-Parametric Models}
We also deploy an RBF-RLS (Recursive Least Squares with 7 Gaussian basis functions) and a Nadaraya-Watson spatial-temporal kernel regression to handle highly non-linear, localized relationships.

\subsection{Portfolio Construction}
Using the outputs from these models, we construct portfolios using mean-variance optimization (MVO) subject to Zero-Beta and Zero-Investment constraints:
\[
\max_w \left( w^T \mu_t - \frac{\gamma}{2} w^T \Sigma_t w \right),
\]
where $\gamma = 200$. Our results demonstrate that the dynamic models (particularly the Kalman polynomial and RBF-RLS) generate superior risk-adjusted returns by dynamically sizing positions based on conditional predictive reliability.

%%Fifth Section
\section{Ethical Implications and Conclusions}

The use of advanced predictive models in finance raises important ethical considerations. If predictive models like these are deployed at scale by large institutional actors, they could induce unintended feedback loops or "flash crashes" due to crowded trades. Additionally, unequal access to high-frequency data and advanced computational tools exacerbates the divide between institutional and retail investors. 

Care must be taken to ensure transparency and to avoid over-reliance on models that may fail under changing market conditions. By maintaining strict risk limits via MVO constraints, we mitigate the risk of catastrophic capital loss. In conclusion, our study of a dynamic conditional Information Coefficient demonstrates that predictive power varies significantly across signal strength and time, and modeling this structure improves portfolio performance.


%%%%%%%%%%%%%%%%%%%%%%%%%%%%%%%%%%%%%
%% Bibliography below
%%%%%%%%%%%%%%%%%%%%%%%%%%%%%%%%%%%%%
\FloatBarrier % Keep the figures from being put after the bibliography
\newpage

\begin{thebibliography}{99}
\bibitem{Vandermeersch} Grinold, R. C., and Kahn, R. N. \emph{Active Portfolio Management: A Quantitative Approach for Producing Superior Returns and Controlling Risk.} McGraw-Hill, 1999.
\end{thebibliography}

\end{document}